\section{Appendix H: Symbol Dynamics \texorpdfstring{$\zeta$}{zeta} Generating Function and Arithmetic Structure Interface (Extension, Non-Main Chain)}
\label{app:zeta-interface}

This appendix provides a rigorous but limited interface: explains ``how prime factorization structure appears in period-orbit counting / $\zeta$ generating function inversion formula,'' without claiming to provide prime distribution theorems or prime-generation algorithms.

\subsection{Artin--Mazur \texorpdfstring{$\zeta$}{zeta} Function of Symbol Dynamics}

For a finite-state shift (SFT), period-orbit counting $\card{\mathrm{Fix}(T^n)}$ defines the $\zeta$ generating function
\[
\zeta_T(z):=\exp\Bigl(\sum_{n\ge 1}\frac{\card{\mathrm{Fix}(T^n)}}{n}z^n\Bigr).
\]
For adjacency matrix $A$, there is a closed form
\[
\zeta_T(z)=\frac{1}{\det(I-zA)}.
\]
For the golden-mean shift (forbidden word $11$), the adjacency matrix is
\[
A=\begin{pmatrix}1&1\\1&0\end{pmatrix},
\quad\Rightarrow\quad
\det(I-zA)=1-z-z^2,
\]
hence
\[
\zeta_T(z)=\frac{1}{1-z-z^2}.
\]

\subsection{Primitive Period Orbits and Möbius Inversion}

Let $P(n)$ be the number of length-$n$ primitive period orbits. The classical relation is
\[
\card{\mathrm{Fix}(T^n)}=\sum_{d\mid n} d\,P(d).
\]
By Möbius inversion,
\[
P(n)=\frac{1}{n}\sum_{d\mid n}\mu(d)\,\card{\mathrm{Fix}(T^{n/d})},
\]
where $\mu$ is the Möbius function. This formula shows: prime factorization structure enters through divisor lattice and Möbius function in the ``primitive--non-primitive'' decomposition. This provides a verifiable arithmetic interface: once period-orbit counts are known, primitive-orbit count follows from inversion containing prime factor information.

\subsection{Relation to Main Text Protocol (Interpretation Layer)}

The main text's ``periodic approximation'' produces finite periodic structures at the parameter-approximation layer; if we further encode these periodic structures as SFT or finite approximations, the $\zeta$ generating function and Möbius inversion become standard tools for ``extracting primitive structure from period counts.'' Interpreting the extracted primitive structure as ``prime generation'' requires additional semantic mapping and external arithmetic object selection, belonging to the interpretation layer, not entering the closed-layer theorem chain.
